% ECONOMICS_PROBLEM: {"id": "D63-59", "title": "EF1 and Pareto optimality for additive mixed manna", "jel_code": "D63", "source_id": "OP-0190"}
% FIRST_POSTED: 2026-10-09
% ATTEMPT_STATUS: unresolved
% ENTRY_KIND: research_attempt
% !TeX program = pdflatex
\documentclass[11pt,a4paper]{article}
\usepackage[T1]{fontenc}
\usepackage[utf8]{inputenc}
\usepackage{lmodern}
\usepackage[margin=24mm]{geometry}
\usepackage{amsmath,amssymb,amsthm,mathtools}
\usepackage{booktabs,longtable,array,enumitem,xurl,hyperref,listings,microtype}
\hypersetup{colorlinks=true,linkcolor=blue,urlcolor=blue,citecolor=blue,pdftitle={Research proof audit}}
\setlength{\parindent}{0pt}
\setlength{\parskip}{5pt}
\setlength{\emergencystretch}{3em}
\setlist{nosep,leftmargin=2em}
\newtheorem{theorem}{Theorem}[section]
\newtheorem{lemma}[theorem]{Lemma}
\newtheorem{proposition}[theorem]{Proposition}
\newtheorem{corollary}[theorem]{Corollary}
\theoremstyle{remark}\newtheorem{remark}[theorem]{Remark}
\newcommand{\R}{\mathbb R}
\newcommand{\N}{\mathbb N}
\newcommand{\E}{\mathbb E}
\newcommand{\Var}{\operatorname{Var}}
\newcommand{\ind}{\mathbf 1}
\newcommand{\EF}{\mathrm{EF}}
\newcommand{\EFM}{\mathrm{EFM}}
\newcommand{\PO}{\mathrm{PO}}
\newcommand{\MMS}{\mathrm{MMS}}
\DeclareMathOperator*{\argmax}{arg\,max}
\DeclareMathOperator*{\argmin}{arg\,min}
\lstset{basicstyle=\ttfamily\footnotesize,breaklines=true,columns=fullflexible,keepspaces=true,frame=single}
\begin{document}
\begin{center}
{\LARGE\bfseries EF1 and Pareto optimality for additive mixed manna}\par\medskip
{\large D63-59.tex}\par
Research attempt and adversarial audit --- 9 October 2026
\end{center}
\label{problem:OP-0190}
\label{game:GT-060}
\textbf{Tools.} Web browsing: YES. Computation: YES (Python, exact integer/rational checks, and explicitly identified numerical diagnostics).
\textbf{Status convention.} A full proof of a restricted theorem does not resolve the full input. Literature results are marked as imported; no novelty or exhaustive current-literature certification is claimed. Computational searches certify only their stated finite domains.

\section{FORMAL RESTATEMENT}
For every integer $n\ge3$, finite item set $M$, and real matrix $(u_{ig})$, let $u_i(S)=\sum_{g\in S}u_{ig}$. Items may have different signs for different agents. A complete allocation $A=(A_i)_{i=1}^n$ is an ordered partition, with empty parts allowed. It is EF1 if, for every ordered pair $i,j$, either $u_i(A_i)\ge u_i(A_j)$ or
\[
\exists g\in A_i\cup A_j:\quad u_i(A_i\setminus\{g\})\ge u_i(A_j\setminus\{g\}).
\]
It is Pareto optimal (PO) if there is no complete partition $B$ for which every $u_i(B_i)\ge u_i(A_i)$ and at least one inequality is strict. The question is
\[
\forall n\ge3\ \forall M\text{ finite}\ \forall u\in\R^{n\times |M|}\quad
\exists A\text{ complete, EF1 and PO}.
\]
This is integral PO; fractional PO is stronger. No subsidies, copies, fractional items, weights or normalization are part of this question. For an empty set of candidate deletions, the EF1 condition requires no envy. Arbitrary nonadditive valuations are not covered by the quantifier on $u$.

\section{QUICK LITERATURE/CONTEXT CHECK}
The supplied survey identifies this mixed-manna fairness/efficiency question \cite{survey}. The accessible January 2026 version of Barman and Verma \cite{bv} proves PO with \emph{introspective} EF1 and explicitly distinguishes that property from the requested either-side deletion EF1. Its abstract still poses ordinary EF1 plus PO for additive mixed manna as unresolved. An introspective witness may add an item from outside the envied bundle, so it cannot silently be used as the witness required here. We do not infer an exhaustive October 2026 status from that source alone.

\section{ATTACK PLAN}
This is an existence problem at the intersection of discrete fair division and Pareto efficiency. The proof track searches for a potential minimizer that preserves a fixed total utility in a common-value subdomain. The disproof track checks small signed instances and tries to break the tempting ``find EF1, then Pareto-improve'' argument.

\begin{center}\begin{tabular}{p{.26\linewidth}p{.65\linewidth}}\toprule Tool&Role\\\midrule
Exact finite enumeration&Checks the smallest genuinely mixed domains.\\
Squared-load potential&Balances a common signed valuation without requiring nonnegative loads.\\
Conservation of total load&Certifies Pareto optimality in a proportional-value subdomain.\\
Itemwise welfare maximization&Handles privately beneficial goods.\\
Single-deletion gap formula&Makes EF1 verification exact for additive utilities.\\
Pareto dominance tests&Separates efficiency from fairness along improvement paths.\\
Weighted-welfare/market methods&Possible general route, but a suitable EF1-preserving construction is not supplied.\\\bottomrule\end{tabular}\end{center}

\section{WORK}
\begin{lemma}[An exact EF1 test]\label{lem:ef1test}
For additive utilities and a fixed pair $i,j$, EF1 is equivalent to
\[
u_i(A_j)-u_i(A_i)\le
\max\bigl(\{0\}\cup\{u_{ig}:g\in A_j\}\cup\{-u_{ig}:g\in A_i\}\bigr).
\]
\end{lemma}
\begin{proof}
A nonpositive gap is no envy. Removing $g$ from $A_j$ lowers the right bundle's value by $u_{ig}$; removing $g$ from $A_i$ raises the left value by $-u_{ig}$. Each permitted deletion therefore repairs the gap precisely when its corresponding quantity is at least the gap. Taking the maximum is exactly the disjunction in the definition. Including zero handles empty bundles.
\end{proof}

\begin{theorem}[Proportional signed preferences]\label{thm:proportional}
For every $n\ge1$, if $u_{ig}=\lambda_i a_g$ for real $a_g$ and strictly positive $\lambda_i$, there is a complete EF1 and PO allocation. In fact every complete allocation is PO, and any minimizer of
\[
\Phi(A)=\sum_{i=1}^n\left(\sum_{g\in A_i}a_g\right)^2
\]
is EF1.
\end{theorem}
\begin{proof}
Write $\ell_i=\sum_{g\in A_i}a_g$. A Pareto improvement would weakly raise every $\ell_i$ and strictly raise at least one, because $\lambda_i>0$. That is impossible since $\sum_i\ell_i=\sum_ga_g$ for every complete allocation. Thus every allocation is PO.

A potential minimizer exists because there are finitely many partitions. Suppose it is not EF1. For some $i,j$, let $D=\ell_j-\ell_i>0$. If $A_j$ contains a positive item of value at least $D$, deleting it repairs this envy; if $A_i$ contains a negative item of absolute value at least $D$, deleting it also repairs the envy. Therefore EF1 failure implies every positive item in $A_j$ has value in $(0,D)$ and every negative item in $A_i$ has absolute value in $(0,D)$. At least one such item exists: otherwise $\ell_j\le0$ and $\ell_i\ge0$, contradicting $D>0$.

Move a positive item of value $r\in(0,D)$ from $A_j$ to $A_i$, or a negative item of value $-r$ from $A_i$ to $A_j$. In either case the two loads become $\ell_i+r$ and $\ell_j-r$. The change in the potential is
\[
(\ell_i+r)^2+(\ell_j-r)^2-\ell_i^2-\ell_j^2=2r^2-2Dr<0,
\]
a contradiction. Since multiplying a pairwise comparison by $\lambda_i>0$ preserves its direction, this proves EF1 for the original utilities.
\end{proof}
The proof establishes existence by a finite minimizer; it does not claim polynomial-time optimization of that potential.

\begin{theorem}[Privately beneficial items]\label{thm:private}
Suppose every item $g$ has a designated agent $h(g)$ with $u_{h(g),g}\ge0$, while $u_{ig}\le0$ for every $i\ne h(g)$. Assigning $g$ to $h(g)$ yields an exactly envy-free and PO allocation.
\end{theorem}
\begin{proof}
Every agent values its own assigned items nonnegatively, so its own utility is at least zero. Every item assigned to somebody else is nonpositive to that agent, so its value for another bundle is at most zero. Hence the allocation is envy-free. The assignment maximizes $\sum_i u_i(A_i)$ item by item, since $h(g)$ has a weakly largest value for $g$. A Pareto improvement would strictly increase this sum, contradicting maximality.
\end{proof}

\begin{proposition}[Pareto improvement need not preserve EF1]\label{prop:badroute}
There are additive nonnegative utilities and complete allocations $A,B$ such that $A$ is EF1, $B$ Pareto dominates $A$, and $B$ is not EF1.
\end{proposition}
\begin{proof}
Use three agents and items $a,b,c$ with value matrix
\[
\begin{array}{c|rrr}&a&b&c\\\hline
1&1&0&1\\2&2&1&2\\3&0&0&0
\end{array}.
\]
Let $A=(\{a\},\{b\},\{c\})$. Every compared bundle has at most one item, so deleting that item leaves zero, no larger than any agent's nonnegative own utility. Hence $A$ is EF1. Set $B=(\{a,c\},\{b\},\varnothing)$. Own utilities change from $(1,1,0)$ to $(2,1,0)$, a Pareto improvement. Agent 2 values the first bundle at 4 and its own at 1. Deleting either item from the first bundle leaves value 2; deleting its own item leaves zero. Neither deletion repairs its envy. Thus $B$ is not EF1.
\end{proof}
This is a counterexample to a proposed \emph{proof step}, not to the existence statement. In particular, the report does not relabel it as a disproof of EF1 plus PO.

\section{VERIFICATION}
The accompanying program exhaustively checked all $3^9=19{,}683$ utility matrices with $n=m=3$ and entries in $\{-1,0,1\}$. For each it checked all $3^3=27$ complete allocations, using Lemma~\ref{lem:ef1test} and pairwise Pareto dominance against all other allocations. Every profile had an EF1-and-PO allocation; the smallest number of such allocations in a profile was one. All comparisons were exact integers.
\begin{lstlisting}
for u in {-1,0,1}^(3 by 3):
    allocations = all 27 complete partitions
    for A in allocations:
        fair[A] = the single-deletion gap test holds for all i,j
        PO[A] = no allocation B weakly improves everyone and strictly improves one
    assert some A satisfies fair[A] and PO[A]
\end{lstlisting}
Zero values, empty bundles, positive and negative total loads, and negative items with agent-dependent signs were included. The proportional theorem permits arbitrary real magnitudes; its use of $\lambda_i>0$ is essential. If some $\lambda_i=0$, the assertion that every allocation is PO need not follow from conserved loads and is not claimed. Finite examples with ternary values cannot certify arbitrary real utilities or larger item sets. The separately supplied Pareto-step counterexample is consistent with the positive pure-goods existence results.

\section{FINAL}
\textbf{LABEL: UNRESOLVED.}
The general heterogeneous additive mixed-manna statement is not proved or disproved. Theorems~\ref{thm:proportional} and \ref{thm:private} are full proofs of explicit subdomains. Proposition~\ref{prop:badroute} rigorously invalidates a natural shortcut.

\textbf{Exact first gap.} For arbitrary signed heterogeneous utility matrices, no selection of an EF1 allocation is proved to lie on the Pareto frontier, and no EF1-preserving Pareto-improvement operation is established.

\textbf{Three next moves.} (1) Seek a weighted-welfare maximizing allocation whose every excessive envy gap admits an efficiency-preserving alternating transfer. (2) Extend the signed potential proof to two nonproportional utility types, locating the first failed conservation identity. (3) Search sparse, rational four-agent profiles with uneven positive and negative magnitudes, retaining exact Pareto/EF1 certificates rather than relying on random dense instances.

\textbf{Minimal hard counterexample, if one exists.} It must escape the two proved subdomains; heterogeneity in signs or relative magnitudes is necessary for defeating the common-load proof. The finite test rules out only the stated three-agent ternary three-item domain, not all ternary markets.

\section{Completion Estimate}
18\%

Two broad structured subcases and an exact small-domain test are complete, but the arbitrary heterogeneous signed-utility existence question is not resolved.

This is a subjective estimate of the fraction of the \emph{whole requested mathematical scope} addressed by this report, including the named variants. It is not a probability of correctness, an objectively measured fraction of a proof, or an estimate that the remaining work takes the complementary fraction of the time. Only the eleven supplied files are assessed; no missing rank-1--200 statements are inferred.

\begin{thebibliography}{99}
\bibitem{survey} S. Liu, X. Lu, M. Suzuki and T. Walsh. \emph{Mixed Fair Division: A Survey}. arXiv:2306.09564v4, Section 4.1, Open Question 2. \url{https://arxiv.org/html/2306.09564v4}.
\bibitem{bv} S. Barman and P. Verma. \emph{Introspectively Envy-Free and Efficient Allocation of Indivisible Mixed Manna}. Accessible arXiv HTML version 2509.18673v3, 2 January 2026, abstract and Definition 1. \url{https://arxiv.org/html/2509.18673}.
\end{thebibliography}
\end{document}
